\section{Алгоритмы сортировки}

\subsection{Постановка задачи сортировки}

\textbf{Сортировкой} называется задача упорядочения элементов некоторой
последовательности в соответствии с заданным отношением порядка.

Пусть дана последовательность
\[
A = (a_1,a_2,\ldots,a_n).
\]
Требуется построить перестановку элементов
\[
A' = (a'_1,a'_2,\ldots,a'_n),
\]
такую, что
\[
a'_1 \leq a'_2 \leq \ldots \leq a'_n.
\]

При этом $A'$ должна содержать те же элементы, что и $A$, то есть отличаться
от исходной последовательности только перестановкой элементов.

На практике сортируются не просто числа, а \textbf{записи}, содержащие
некоторый ключ:
\[
a_i = (\text{key}_i,\text{data}_i).
\]
Порядок элементов определяется значениями ключей.

Сортировка является одной из фундаментальных алгоритмических задач.
Она используется для:
\begin{itemize}
	\item ускорения поиска;
	\item группировки одинаковых элементов;
	\item построения индексов;
	\item обработки баз данных;
	\item подготовки данных к другим алгоритмам;
	\item нахождения медиан, квантилей и порядковых статистик.
\end{itemize}

Алгоритмы сортировки существенно различаются по вычислительной сложности,
объёму дополнительной памяти и требованиям к исходным данным.


\subsection{Устойчивость и другие свойства алгоритмов}

При выборе алгоритма сортировки важна не только его асимптотическая
сложность.

\textbf{Устойчивая сортировка} сохраняет взаимный порядок элементов
с одинаковыми ключами.

Например, пусть записи имеют вид
\[
(2,A),\quad (1,B),\quad (2,C).
\]
После устойчивой сортировки по первому полю получим
\[
(1,B),\quad (2,A),\quad (2,C).
\]
Записи $(2,A)$ и $(2,C)$ сохранили исходный порядок.

Устойчивость особенно важна при последовательной сортировке по нескольким
ключам. Например, если записи сначала отсортировать по имени, а затем
устойчиво по фамилии, то внутри одинаковых фамилий сохранится сортировка
по имени.

Основные характеристики алгоритмов сортировки:

\begin{itemize}
	\item \textbf{временная сложность} --- число операций в зависимости от $n$;
	\item \textbf{пространственная сложность} --- объём дополнительной памяти;
	\item \textbf{устойчивость};
	\item \textbf{сортировка на месте} (\textit{in-place}) --- использование
	$O(1)$ дополнительной памяти, не считая стека рекурсии;
	\item \textbf{адаптивность} --- способность работать быстрее на частично
	упорядоченных данных;
	\item \textbf{онлайновость} --- возможность обрабатывать элементы по мере
	их поступления;
	\item использование или отсутствие операций сравнения.
\end{itemize}

Сортировки обычно разделяют на две большие группы:

\begin{enumerate}
	\item \textbf{сортировки сравнениями}, в которых информация о порядке
	элементов получается с помощью операций вида
	\[
	a_i < a_j;
	\]
	\item \textbf{сортировки без сравнений}, использующие дополнительные
	свойства ключей.
\end{enumerate}


\subsection{Сортировки вставками, выбором и обменами}

\subsubsection*{Сортировка вставками}

Идея \textbf{сортировки вставками} (\textit{Insertion Sort}) состоит в том,
что на каждом шаге очередной элемент вставляется в правильное место уже
отсортированной части массива.

Пусть элементы
\[
a_1,\ldots,a_{i-1}
\]
уже отсортированы. Элемент $a_i$ сравнивается с предыдущими элементами и
сдвигается влево до тех пор, пока не займёт требуемое положение.

Для каждого элемента может потребоваться просмотр почти всей
предшествующей части массива.

В худшем случае:
\[
T(n)=1+2+\ldots+(n-1)
=\frac{n(n-1)}{2}
=\Theta(n^2).
\]

В лучшем случае, когда массив уже отсортирован,
\[
T(n)=\Theta(n).
\]

Свойства:
\begin{itemize}
	\item худшая сложность --- $O(n^2)$;
	\item средняя сложность --- $O(n^2)$;
	\item лучшая сложность --- $O(n)$;
	\item дополнительная память --- $O(1)$;
	\item сортировка устойчива;
	\item алгоритм адаптивен;
	\item эффективен для небольших или почти отсортированных массивов.
\end{itemize}


\subsubsection*{Сортировка выбором}

В \textbf{сортировке выбором} (\textit{Selection Sort}) на каждом шаге
из ещё не отсортированной части массива выбирается минимальный элемент и
переставляется в начало этой части.

На первом шаге ищется минимум среди $n$ элементов, затем среди $n-1$ и так
далее.

Число сравнений:
\[
(n-1)+(n-2)+\ldots+1
=\frac{n(n-1)}{2}.
\]

Следовательно,
\[
T(n)=\Theta(n^2)
\]
при любом расположении элементов.

При этом число перестановок равно лишь $O(n)$.

Свойства:
\begin{itemize}
	\item сложность во всех случаях --- $\Theta(n^2)$;
	\item дополнительная память --- $O(1)$;
	\item обычная реализация неустойчива;
	\item число обменов мало.
\end{itemize}

Алгоритм может быть полезен, когда операция перемещения элемента значительно
дороже его сравнения.


\subsubsection*{Сортировка обменами}

Наиболее известным алгоритмом этого класса является
\textbf{пузырьковая сортировка} (\textit{Bubble Sort}).

Последовательно сравниваются соседние элементы:
\[
a_i,\;a_{i+1}.
\]
Если
\[
a_i>a_{i+1},
\]
то они меняются местами.

После одного полного прохода максимальный элемент оказывается в конце
массива. После следующего прохода на своё место перемещается второй по
величине элемент и т.д.

В худшем и среднем случаях:
\[
T(n)=\Theta(n^2).
\]

При использовании признака того, что за проход не произошло ни одного
обмена, для уже отсортированного массива получается
\[
T(n)=\Theta(n).
\]

Свойства:
\begin{itemize}
	\item простота реализации;
	\item устойчивость;
	\item $O(1)$ дополнительной памяти;
	\item низкая практическая эффективность при больших $n$.
\end{itemize}

Поэтому пузырьковая сортировка имеет преимущественно учебное значение.


\subsection{Сортировка слиянием}

\textbf{Сортировка слиянием} (\textit{Merge Sort}) основана на принципе
``разделяй и властвуй''.

Алгоритм состоит из трёх этапов:

\begin{enumerate}
	\item разбить массив примерно пополам;
	\item рекурсивно отсортировать обе половины;
	\item слить две отсортированные последовательности.
\end{enumerate}

Пусть имеются две отсортированные последовательности
\[
A=(a_1,\ldots,a_k), \qquad
B=(b_1,\ldots,b_m).
\]

При слиянии сравниваются первые ещё не использованные элементы обеих
последовательностей. Меньший из них добавляется в результирующий массив.

Слияние выполняется за
\[
\Theta(k+m).
\]

Для всего алгоритма получается рекуррентное соотношение
\[
T(n)=2T\left(\frac{n}{2}\right)+\Theta(n).
\]

Отсюда
\[
T(n)=\Theta(n\log n).
\]

Это справедливо для лучшего, среднего и худшего случаев.

Свойства Merge Sort:
\begin{itemize}
	\item время работы --- $\Theta(n\log n)$;
	\item классическая реализация требует $O(n)$ дополнительной памяти;
	\item алгоритм устойчив;
	\item хорошо подходит для связных списков;
	\item эффективен для внешней сортировки;
	\item хорошо допускает распараллеливание.
\end{itemize}

В отличие от быстрой сортировки, сортировка слиянием имеет гарантированную
сложность $O(n\log n)$.

\subsection{Быстрая сортировка}

\textbf{Быстрая сортировка} (\textit{Quick Sort}) также основана на принципе
``разделяй и властвуй''.

Выбирается некоторый элемент массива, называемый
\textbf{опорным элементом} (\textit{pivot}).

Затем массив разделяется так, чтобы элементы с меньшими ключами находились
с одной стороны от опорного элемента, а с большими --- с другой.

После разделения обе части сортируются рекурсивно.

Если разделение массива происходит примерно пополам, то
\[
T(n)=2T\left(\frac{n}{2}\right)+\Theta(n),
\]
откуда
\[
T(n)=\Theta(n\log n).
\]

Однако при неудачном выборе опорного элемента одна часть может содержать
$n-1$ элементов, а другая --- ни одного. Тогда
\[
T(n)=T(n-1)+\Theta(n),
\]
и
\[
T(n)=\Theta(n^2).
\]

Таким образом:
\[
T_{\text{avg}}(n)=\Theta(n\log n),
\qquad
T_{\text{worst}}(n)=\Theta(n^2).
\]

Для уменьшения вероятности плохого разбиения применяют:
\begin{itemize}
	\item случайный выбор опорного элемента;
	\item выбор среднего из первого, среднего и последнего элементов;
	\item другие эвристики выбора pivot.
\end{itemize}

Свойства Quick Sort:
\begin{itemize}
	\item среднее время --- $\Theta(n\log n)$;
	\item худшее время --- $\Theta(n^2)$;
	\item обычно работает на месте;
	\item средний размер стека рекурсии --- $O(\log n)$;
	\item классический вариант неустойчив;
	\item обладает хорошей локальностью обращений к памяти;
	\item на практике часто является одним из наиболее быстрых алгоритмов
	сортировки массивов в оперативной памяти.
\end{itemize}

Несмотря на квадратичный худший случай, Quick Sort часто превосходит
Merge Sort и Heap Sort на массивах за счёт малой константы и эффективного
использования кэш-памяти.


\subsection{Пирамидальная сортировка}

\textbf{Пирамидальная сортировка} (\textit{Heap Sort}) использует структуру
данных \textbf{двоичная куча}.

Максимальной кучей называется двоичное дерево, для каждой вершины которого
выполнено
\[
a_{\text{parent}} \geq a_{\text{child}}.
\]

Поэтому максимальный элемент всегда расположен в корне.

Если куча хранится в массиве, то для вершины с индексом $i$ индексы потомков
можно определить арифметически. Например, при нумерации с единицы:
\[
\operatorname{left}(i)=2i,
\qquad
\operatorname{right}(i)=2i+1.
\]

Сортировка состоит из двух этапов:

\begin{enumerate}
	\item из исходного массива строится максимальная куча;
	\item корневой элемент последовательно меняется с последним элементом
	ещё не отсортированной части, после чего восстанавливается свойство кучи.
\end{enumerate}

Кучу из $n$ элементов можно построить за
\[
O(n).
\]

Извлечение максимального элемента и восстановление кучи требует
\[
O(\log n).
\]

Такая операция выполняется $n$ раз, поэтому
\[
T(n)=O(n\log n).
\]

Причём эта оценка справедлива и в худшем случае.

Свойства:
\begin{itemize}
	\item гарантированное время --- $O(n\log n)$;
	\item дополнительная память --- $O(1)$;
	\item алгоритм неустойчив;
	\item не требует дополнительного массива;
	\item обычно медленнее Quick Sort на практике из-за худшей локальности
	памяти.
\end{itemize}


\subsection{Сортировки без сравнений}

Нижняя оценка $O(n\log n)$ относится только к алгоритмам,
основанным на сравнениях.

Если о ключах известна дополнительная информация, возможна сортировка за
линейное время относительно числа элементов.

\subsubsection*{Сортировка подсчётом}

Пусть ключи являются целыми числами из диапазона
\[
0,\ldots,k-1.
\]

Создаётся массив счётчиков
\[
C[0],\ldots,C[k-1],
\]
где $C[i]$ содержит число элементов с ключом $i$.

После подсчёта элементов массив можно восстановить в отсортированном виде.

Сложность:
\[
T(n,k)=O(n+k),
\]
память:
\[
O(k).
\]

При подходящей реализации алгоритм является устойчивым.

Если $k=O(n)$, время работы линейно:
\[
T(n)=O(n).
\]


\subsubsection*{Поразрядная сортировка}

\textbf{Radix Sort} сортирует ключи последовательно по отдельным разрядам.

Например, для десятичных чисел можно сначала выполнить устойчивую сортировку
по единицам, затем по десяткам, сотням и т.д.

Если число разрядов равно $d$, а сортировка одного разряда требует
$O(n+k)$ операций, то
\[
T(n)=O(d(n+k)).
\]

В качестве внутреннего алгоритма часто используется устойчивая сортировка
подсчётом.


\subsubsection*{Карманная сортировка}

В \textbf{Bucket Sort} диапазон возможных значений делится на несколько
интервалов --- ``карманов''. Элементы распределяются по карманам, каждый
карман сортируется отдельно, после чего результаты объединяются.

При подходящем распределении данных среднее время может составлять
\[
O(n).
\]

Однако эффективность сильно зависит от распределения входных значений.


\subsection{Оценка временной и пространственной сложности}

Основные алгоритмы можно сравнить следующим образом:

\begin{center}
	\begin{tabular}{|l|c|c|c|c|}
		\hline
		Алгоритм &
		Лучший случай &
		Средний случай &
		Худший случай &
		Память \\ \hline
		
		Вставки &
		$O(n)$ &
		$O(n^2)$ &
		$O(n^2)$ &
		$O(1)$ \\ \hline
		
		Выбор &
		$O(n^2)$ &
		$O(n^2)$ &
		$O(n^2)$ &
		$O(1)$ \\ \hline
		
		Пузырьковая &
		$O(n)$ &
		$O(n^2)$ &
		$O(n^2)$ &
		$O(1)$ \\ \hline
		
		Слияние &
		$O(n\log n)$ &
		$O(n\log n)$ &
		$O(n\log n)$ &
		$O(n)$ \\ \hline
		
		Быстрая &
		$O(n\log n)$ &
		$O(n\log n)$ &
		$O(n^2)$ &
		$O(\log n)$* \\ \hline
		
		Пирамидальная &
		$O(n\log n)$ &
		$O(n\log n)$ &
		$O(n\log n)$ &
		$O(1)$ \\ \hline
		
		Подсчётом &
		$O(n+k)$ &
		$O(n+k)$ &
		$O(n+k)$ &
		$O(k)$ \\ \hline
	\end{tabular}
\end{center}

Здесь $k$ --- размер диапазона ключей.

Для Quick Sort указана типичная глубина стека рекурсии. В худшем случае
она может достигать $O(n)$, если специально не ограничивать рекурсию.

Асимптотическая сложность является важной, но не единственной характеристикой.
Например, два алгоритма с одинаковой сложностью
\[
O(n\log n)
\]
могут существенно различаться на практике из-за:
\begin{itemize}
	\item числа операций сравнения;
	\item числа перемещений элементов;
	\item характера обращений к памяти;
	\item использования кэш-памяти;
	\item стоимости выделения дополнительной памяти.
\end{itemize}


\subsection{Нижняя оценка для сортировок сравнениями}

Рассмотрим произвольный алгоритм сортировки, который получает информацию
о взаимном порядке элементов только посредством сравнений.

Работу такого алгоритма можно представить в виде
\textbf{дерева решений}.

Каждая внутренняя вершина соответствует сравнению двух элементов.
В зависимости от результата сравнения алгоритм переходит в одну из ветвей.

Для $n$ различных элементов существует
\[
n!
\]
возможных перестановок.

Чтобы алгоритм мог отличить каждую перестановку от всех остальных,
дерево решений должно иметь не менее $n!$ листьев.

Если высота бинарного дерева равна $h$, то оно имеет не более
\[
2^h
\]
листьев.

Следовательно,
\[
2^h \geq n!,
\]
откуда
\[
h \geq \log_2(n!).
\]

По формуле Стирлинга
\[
n! \sim \sqrt{2\pi n}
\left(\frac{n}{e}\right)^n,
\]
поэтому
\[
\log_2(n!)=\Theta(n\log n).
\]

Таким образом, для любого алгоритма сортировки сравнениями существует
вход, на котором необходимо выполнить не менее
\[
\Omega(n\log n)
\]
сравнений.

Следовательно, асимптотически оптимальные сортировки сравнениями имеют
сложность
\[
\Theta(n\log n).
\]

К ним относятся, например, Merge Sort и Heap Sort.

Это ограничение не противоречит существованию сортировки подсчётом или
Radix Sort с линейной сложностью, поскольку они используют дополнительную
информацию о структуре ключей и не являются чистыми сортировками
сравнениями.


\subsection{Внешняя сортировка и выбор алгоритма}

Если все данные помещаются в оперативную память, говорят о
\textbf{внутренней сортировке}.

Если объём данных превышает доступную оперативную память и существенная
часть данных хранится на внешнем устройстве, применяется
\textbf{внешняя сортировка}.

В этом случае основной стоимостью становятся не арифметические операции,
а операции ввода-вывода.

Классическим методом является
\textbf{внешняя сортировка слиянием}.

Она выполняется следующим образом:

\begin{enumerate}
	\item из внешней памяти считывается блок данных, помещающийся в ОЗУ;
	\item блок сортируется внутренним алгоритмом;
	\item отсортированный блок записывается на диск; такой блок называется
	\textbf{серией} или \textbf{run};
	\item полученные серии объединяются многопутевым слиянием;
	\item процесс повторяется до получения одной отсортированной
	последовательности.
\end{enumerate}

Слияние особенно удобно для внешней памяти, поскольку данные читаются
последовательно большими блоками.

При выборе алгоритма сортировки следует учитывать свойства задачи.

Для небольших массивов хорошо подходит сортировка вставками.

Для произвольных массивов в оперативной памяти на практике часто
эффективна Quick Sort или гибридные алгоритмы на её основе.

Если требуется гарантированное
\[
O(n\log n)
\]
и минимальная дополнительная память, может использоваться Heap Sort.

Если важна устойчивость и допустимо выделение дополнительной памяти,
подходит Merge Sort.

Если ключи --- целые числа из небольшого диапазона, эффективна сортировка
подсчётом.

Для составных целочисленных или строковых ключей может применяться
Radix Sort.

Для данных, не помещающихся в оперативной памяти, естественным выбором
является внешняя сортировка слиянием.

Таким образом, универсально ``лучшего'' алгоритма сортировки не существует:
выбор определяется размером данных, типом ключей, требованием устойчивости,
ограничениями памяти и архитектурой вычислительной системы.

\newpage